\section{The complete theta source and logarithmic resolvent}\label{sec:source}

\subsection{Reflection, the endpoint and the canonical normalization}

Let
\[
 \vartheta(t)=\sum_{n\in\mathbb Z}\ee^{-\pi n^2t},\qquad
 p(u)=\frac{\ee^u}{2}\{\vartheta(\ee^{4u})-1\}
     =\ee^u\sum_{n\ge1}\ee^{-\pi n^2\ee^{4u}}.
\]
Jacobi's transformation
\(\vartheta(t)=t^{-1/2}\vartheta(1/t)\)
\cite{DLMFThetaModular} gives
\begin{equation}\label{eq:primitive-reflection}
 p(-u)-p(u)=\sinh u,\qquad p'(0)=-\frac12.
\end{equation}
Define
\begin{equation}\label{eq:source}
 \Phi(u)=\frac{p''(u)-p(u)}8
 =\ee^u\sum_{n\ge1}(2r_n(u)^2-3r_n(u))\ee^{-r_n(u)},
 \qquad r_n(u)=\pi n^2\ee^{4u}.
\end{equation}
Applying \(\partial_u^2-1\) to \eqref{eq:primitive-reflection}
proves that \(\Phi\) is even. The complete series is smooth, and
each summand is positive when \(u\ge0\), since \(r_n\ge\pi>3/2\).
Thus \(\Phi>0\) on \(\R\). Reflection is a property of the complete
sum; an evenly reflected finite prefix need not have matching odd
derivatives at zero.

With \(W_\lambda(u)=\ee^{\lambda u^2}\Phi(u)\), put
\begin{equation}\label{eq:H}
 H_\lambda(z)=\frac12\int_\R W_\lambda(u)\ee^{\ii zu}\dd u
            =\int_0^\infty W_\lambda(u)\cos(zu)\dd u.
\end{equation}
There is no division by \(H_\lambda(0)\).
For \(J^{\rm prim}(z)=\int_0^\infty p(u)\cos(zu)\dd u\),
two integrations by parts, using the endpoint in
\eqref{eq:primitive-reflection}, give
\begin{equation}\label{eq:primitive-H}
 8H_0(z)=\frac12-(1+z^2)J^{\rm prim}(z).
\end{equation}
For completeness, splitting the theta Mellin integral at \(t=1\)
and using Jacobi reflection gives
\[
 \xi(s)=\frac12+\frac{s(s-1)}2
 \int_1^\infty \frac{\vartheta(t)-1}{2}
       \{t^{s/2-1}+t^{(1-s)/2-1}\}\dd t.
\]
The integral is entire in \(s\). Substituting \(t=\ee^{4u}\)
and then \(s=(1-\ii z)/2\) turns the integral into
\(8J^{\rm prim}(z)\), while \(4s(s-1)=-(1+z^2)\).
Equation~\eqref{eq:primitive-H} therefore proves
\begin{equation}\label{eq:xi}
 H_0(z)=\frac18\xi(s(z)),\quad
 s(z)=\frac12-\frac{\ii z}2,\quad
 \xi(s)=\frac12s(s-1)\pi^{-s/2}\Gamma(s/2)\zeta(s).
\end{equation}
The reflected argument gives the same field
\cite{DLMFGreenPrime}. The constant \(1/2\) in
\eqref{eq:primitive-H} is indispensable.

\subsection{Complete domination, heat evolution and the product}

Set \(\rho(u)=\pi\ee^{4|u|}\), and recursively define
\[
 P_0(r)=2r^2-3r,\qquad
 P_{j+1}(r)=(1-4r)P_j(r)+4rP_j'(r).
\]
Write \(a(P)\) for the sum of the absolute coefficients of a polynomial,
and put
\[
 C_d(N)=\sum_{n>N}n^{2d}\ee^{-\pi(n^2-1)},\qquad
 B_j=a(P_j)C_{j+2}(0).
\]
Each \(C_d(N)\) is finite and tends to zero with \(N\).

\begin{lemma}[Complete source bound]\label{lem:source}
For each fixed integer \(j\ge0\),
\[
 |\Phi^{(j)}(u)|\le B_j\ee^{|u|}\rho(u)^{j+2}\ee^{-\rho(u)}.
\]
On \(u\ge0\), the differentiated arithmetic tail satisfies the same
bound with \(B_j\) replaced by \(a(P_j)C_{j+2}(N)\).
\end{lemma}
\begin{proof}
Differentiating a summand gives
\(\ee^uP_j(n^2\rho)\ee^{-n^2\rho}\).
Since \(\deg P_j=j+2\) and \(\rho\ge\pi>1\),
\[
 |P_j(n^2\rho)|\le a(P_j)n^{2j+4}\rho^{j+2},\qquad
 \ee^{-n^2\rho}\le\ee^{-\rho}\ee^{-\pi(n^2-1)}.
\]
Summation proves the half-line estimate and differentiated local
convergence. Complete smooth reflection gives the full-line estimate.
\end{proof}

For fixed \(L,b\ge0\) and integers \(k,d\ge0\), the integral
\[
 \int_\R |u|^k\ee^{Lu^2+(b+1)|u|}
             \rho(u)^d\ee^{-\rho(u)}\dd u
\]
is finite: substitution \(t=\pi\ee^{4|u|}\) leaves a product of
powers of \(t,\log(t/\pi)\), and
\(\exp\{(L/16)\log^2(t/\pi)-t\}\).
Consequently \(H\) is jointly entire in \((\lambda,z)\), and all
fixed derivatives on compact parameter sets pass under the integral.
In particular,
\begin{equation}\label{eq:heat}
 \partial_\lambda H_\lambda=-\partial_z^2H_\lambda.
\end{equation}
The opposite clock \(h(\tau,z)=H_{\lambda_0-\tau}(z)\) satisfies
\(h_\tau=h_{zz}\). No uniform assertion in a growing derivative
order follows from these fixed-order bounds.

We retain the canonical strip input
\begin{equation}\label{eq:zero-strip}
 |\Im z_j(\lambda)|\le1\qquad(0\le\lambda\le1/2)
\end{equation}
from the parent development \cite{OGParent,GreenStripProof}.
This is strip control, not an all-real-zero assumption.
The source bound also gives, uniformly in this compact time range,
\(\log\max_{|z|\le r}|H_\lambda(z)|=O(r\log(r+2))\).
One can obtain this bound directly by splitting the source integral
at a constant multiple of \(\log(r+2)\); beyond that point the
double exponential dominates \(r|u|+u^2/2+O(|u|)\).
Positivity gives a common lower bound for \(H_\lambda(0)\).
Jensen's formula on a disk of radius \(2r\) therefore gives a
full origin-centered count
\begin{equation}\label{eq:count}
 N_\lambda(r)\le Cr\log(r+2),\qquad r\ge1.
\end{equation}
The even Hadamard product is consequently
\begin{equation}\label{eq:product}
 H_\lambda(z)=H_\lambda(0)
 \prod_{\zeta\in\mathcal R_\lambda}(1-z^2/\zeta^2).
\end{equation}
Here \(\mathcal R_\lambda\) contains one representative of each
opposite pair, with multiplicity. The paired product converges
locally uniformly because \(\sum|\zeta|^{-2}<\infty\). The order
bound permits at most a linear exponential factor; evenness removes
its linear term, and evaluation at zero fixes the constant.
Thus no unspecified exponential factor remains.

\subsection{The logarithmic field and its denominator}

Let \(\nu_\lambda=\sum_j\delta_{z_j(\lambda)}\), with multiplicities,
and use \(\partial_z=(\partial_x-\ii\partial_y)/2\).
Local factorization \(H=(z-z_j)^m h\), with \(h\ne0\), proves
\begin{equation}\label{eq:distributions}
 \Delta\log|H_\lambda|=2\pi\nu_\lambda,\qquad
 R_\lambda=\frac{H_\lambda'}{H_\lambda}
           =2\partial_z\log|H_\lambda|,\qquad
 \partial_{\bar z}R_\lambda=\pi\nu_\lambda .
\end{equation}
These are distributional identities; \(\log|z|\) and \(1/z\) are
locally integrable, with the stated delta masses obtained by integrating
around a small circle. A multiplicity-\(m\) zero gives residue \(m\).
Differentiating \eqref{eq:product} away from its zeros gives
\[
 R_\lambda(z)=\sum_{\zeta\in\mathcal R_\lambda}
                    \frac{2z}{z^2-\zeta^2}.
\]
Pairing is part of the convergence prescription.

On zero-free space-time patches, logarithmic differentiation of
\eqref{eq:heat} yields
\begin{equation}\label{eq:burgers}
 \partial_\lambda R=-R_{zz}-2RR_z.
\end{equation}
This local equation is not assigned an additional unproved global
meaning at its poles. Likewise, solving only the first or last
distributional equation in \eqref{eq:distributions} leaves a harmonic
or holomorphic summand. The complete field and its product fix it.

Source approximation alone does not control division. If
\(|H-\widehat H|\le\varepsilon\),
\(|H'-\widehat H'|\le\varepsilon_1\), and
\(b=|\widehat H|-\varepsilon>0\), then
\begin{equation}\label{eq:ratio-error}
 \left|R-\frac{\widehat H'}{\widehat H}\right|
 \le\frac{\varepsilon_1}{b}
       +\frac{|\widehat H'|\varepsilon}{b|\widehat H|}.
\end{equation}
This follows by adding and subtracting
\(\widehat H'/H\). The positive denominator margin is a separate
requirement. It fails on a real interval containing a zero, which
motivates the integrated logarithmic method developed later.

